% Part: incompleteness
% Chapter: theories-computability
% Section: extensions-of-q-not-decidable

\documentclass[../../../include/open-logic-section]{subfiles}

\begin{document}

\olfileid[hi]{inc}{tcp}{cqn}

\olsection{$\Th{Q}$ के संगत विस्तार अनिर्णेय हैं}

\begin{explain}
याद करें कि कोई सिद्धांत \emph{संगत} है यदि वह किसी सूत्र~$!A$ के लिए
$!A$ और $\lnot !A$ दोनों को सिद्ध नहीं करता। चूँकि विरोधाभास से कुछ भी
निष्पन्न होता है, असंगत सिद्धांत तुच्छ है: प्रत्येक वाक्य सिद्ध होता है।
स्पष्टतः यदि कोई सिद्धांत $\omega$-संगत है, तो वह संगत है। किंतु संगत होना
दुर्बलतर आवश्यकता है (अर्थात् ऐसे सिद्धांत हैं जो संगत तो हैं, पर
$\omega$-संगत नहीं)। हम \olref[oqn]{thm:oconsis-q} की मान्यता को केवल
संगतता तक दुर्बल करके अधिक प्रबल प्रमेय प्राप्त कर सकते हैं।
\end{explain}

\begin{lem}
कोई ``सार्वत्रिक संगणनीय संबंध'' नहीं है। अर्थात् ऐसा कोई द्वि-स्थानीय
संगणनीय संबंध $R(x,y)$ नहीं है जिसमें यह गुण हो: जब भी $S(y)$ कोई
एक-स्थानीय संगणनीय संबंध हो, तब कोई $k$ ऐसा हो कि प्रत्येक $y$ के लिए
$S(y)$ तभी और केवल तभी सत्य हो जब $R(k,y)$ सत्य हो।
\end{lem}

\begin{proof}
मान लें $R(x,y)$ कोई सार्वत्रिक संगणनीय संबंध है। $S(y)$ को संबंध
$\lnot R(y,y)$ मानें। चूँकि $S(y)$ संगणनीय है, किसी $k$ के लिए $S(y)$,
$R(k,y)$ के समतुल्य है। किंतु तब $S(k)$,~$R(k,k)$ और $\lnot R(k,k)$ दोनों
के समतुल्य है, जो विरोधाभास है।
\end{proof}

\begin{thm}
$\Th{T}$ कोई ऐसा संगत सिद्धांत हो जिसमें $\Th{Q}$ सम्मिलित है। तब
$\Th{T}$ निर्णेय नहीं है।
\end{thm}

\begin{proof}
मान लें $\Th{T}$,~$\Th{Q}$ का कोई संगत, निर्णेय विस्तार है। हम $\Th{T}$
का उपयोग करके सार्वत्रिक संगणनीय संबंध परिभाषित करेंगे और विरोधाभास
प्राप्त करेंगे।

$R(x,y)$ को तभी और केवल तभी लागू मानें जब
\begin{quote}
$x$ किसी सूत्र $!D(u)$ का कूट है, और $\Th{T}$,~$!D(\num y)$ को सिद्ध करता
है।
\end{quote}
हमारी मान्यता के अनुसार $\Th{T}$ निर्णेय है, इसलिए $R$ संगणनीय है। दिखाएँ
कि $R$ सार्वत्रिक है। यदि $S(y)$ कोई भी संगणनीय संबंध है, तो $\Th{Q}$ में
(और इसलिए $\Th{T}$ में) उसका निरूपण किसी सूत्र $!D_S(u)$ द्वारा होता है।
तब प्रत्येक $n$ के लिए हमारे पास
\begin{eqnarray*}
S(\num n) & \lif & T \vdash !D_S(\num n) \\
& \lif & R(\Gn{!D_S(u)}, n)
\end{eqnarray*}
और
\begin{eqnarray*}
\lnot S(\num n) & \lif & T \vdash \lnot !D_S(\num n) \\
& \lif & T \not\vdash !D_S(\num n) \quad \text{(क्योंकि $\Th{T}$ संगत है)} \\
& \lif & \lnot R(\Gn{!D_S(u)}, n).
\end{eqnarray*}
अर्थात् प्रत्येक $y$ के लिए $S(y)$ तभी और केवल तभी सत्य है जब
$R(\Gn{!D_S(u)}, y)$ सत्य है। अतः $R$ सार्वत्रिक है, और हमें वांछित
विरोधाभास मिल गया।
\end{proof}

``सत्य अंकगणित'' को सिद्धांत $\Setabs{!A}{ \Sat{\Nat}{!A}}$ मानें,
अर्थात् अंकगणित की भाषा के उन वाक्यों का समुच्चय जो मानक व्याख्या में सत्य
हैं।

\begin{cor}
सत्य अंकगणित निर्णेय नहीं है।
\end{cor}

%In Section~\olref{undefinability:truth:section} we will state a stronger
%result, due to Tarski.

\end{document}
